\section{ضرب تانسوری: دومدول‌ها و یکریختی‌های طبیعی}

% Source: sources/lectures/20 tensor product 2.pdf | PDF page: 1 | item: definition of an R-S bimodule
\begin{frame}[t]{دومدول}
\label{lec20:bimodule}
\begin{definition}
فرض کنید \(R\) و \(S\) دو حلقه باشند. گروه آبلی \(A\) را یک
\textbf{\(R\)-\(S\) دومدول} \lr{(bimodule)} می‌گوییم هرگاه \(A\) یک
\(R\)-مدول چپ و یک \(S\)-مدول راست باشد و
\[
    (ra)s=r(as)\qquad(a\in A,\ r\in R,\ s\in S).
\]
دومدول \(A\) را با نماد \({}_R A_S\) نمایش می‌دهیم.
\end{definition}
\begin{remark}
داشتن یک ساختار مدول چپ و یک ساختار مدول راست روی یک گروه آبلی، لزوماً آن را دومدول نمی‌کند.
\end{remark}
\end{frame}

% Source: sources/lectures/20 tensor product 2.pdf | PDF page: 1 | item: separate left and transpose-right actions on M_2(R)
\begin{frame}[t]{دو ساختار مدولی روی ماتریس‌ها}
\label{lec20:matrix-actions}
بگذارید
\[
    A=R=S=M_2(\mathbb R).
\]
ضرب معمول ماتریس‌ها، ساختار \(R\)-مدول چپ را می‌دهد:
\[
    Y\cdot X=YX.
\]
ساختار \(S\)-مدول راست را به صورت زیر تعریف می‌کنیم:
\[
    X\cdot Z=Z^tX.
\]
در این صورت
\[
    (X\cdot Z_1)\cdot Z_2
    =Z_2^tZ_1^tX=(Z_1Z_2)^tX=X\cdot(Z_1Z_2).
\]
\end{frame}

% Source: sources/lectures/20 tensor product 2.pdf | PDF page: 1 | item: failure of compatibility of the two matrix actions
\begin{frame}[t]{ناسازگاری دو عمل}
\label{lec20:matrix-incompatibility}
برای دو عملِ تعریف‌شده داریم
\[
    (Y\cdot X)\cdot Z=(YX)\cdot Z=Z^tYX,
\]
ولی
\[
    Y\cdot(X\cdot Z)=Y\cdot(Z^tX)=YZ^tX.
\]
این دو عبارت لزوماً برابر نیستند. پس \(A\) با این دو عمل، یک \(R\)-\(S\) دومدول نیست.
\end{frame}

% Source: sources/lectures/20 tensor product 2.pdf | PDF page: 2 | item: left S action and induced S-linear tensor morphism
\begin{frame}[t]{ساختار مدول چپ روی ضرب تانسوری}
\label{lec20:left-tensor-action}
\begin{theorem}
فرض کنید \(R,S\) دو حلقه و \({}_S A_R\)، \({}_R B\) مدول‌های داده‌شده باشند.
\begin{enumerate}
\item \(A\otimes_R B\) با عمل
\[
    s\cdot(a\otimes b)=(sa)\otimes b
\]
یک \(S\)-مدول چپ است.
\item اگر \(f:A\to A'\) یک همریختی دومدول‌های \(S\)-\(R\) و
\(g:B\to B'\) یک همریختی \(R\)-مدولی باشد، آنگاه
\[
    f\otimes g:A\otimes_R B\longrightarrow A'\otimes_R B'
\]
یک همریختی \(S\)-مدولی است.
\end{enumerate}
\end{theorem}
\end{frame}

% Source: sources/lectures/20 tensor product 2.pdf | PDF page: 2 | item: balanced map alpha_s
\begin{frame}[t]{خوش‌تعریف بودن عمل چپ}
\label{lec20:alpha-balanced}
\textbf{اثبات.} برای \(s\in S\)، تابع زیر را در نظر بگیرید:
\[
    \alpha_s:A\times B\longrightarrow A\otimes_R B,
    \qquad \alpha_s(a,b)=(sa)\otimes b.
\]
این تابع در هر متغیر جمع‌پذیر است؛ برای نمونه
\[
\begin{aligned}
\alpha_s(a_1+a_2,b)
  &=(sa_1+sa_2)\otimes b\\
  &=(sa_1)\otimes b+(sa_2)\otimes b.
\end{aligned}
\]
همچنین، به علت دومدول بودن \(A\)،
\[
    \alpha_s(ar,b)=(s(ar))\otimes b
    =((sa)r)\otimes b=(sa)\otimes rb=\alpha_s(a,rb).
\]
پس \(\alpha_s\) تابعی خطی میانی است.
\end{frame}

% Source: sources/lectures/20 tensor product 2.pdf | PDF page: 3 | item: universal map alpha-bar and left module axioms
\begin{frame}[t]{اصول مدولی برای عمل القایی}
\label{lec20:left-module-axioms}
خاصیت جهانی، همریختی گروهی زیر را به دست می‌دهد:
\[
    \overline{\alpha_s}:A\otimes_R B\longrightarrow A\otimes_R B,
    \qquad a\otimes b\longmapsto (sa)\otimes b.
\]
تعریف می‌کنیم \(s\cdot x=\overline{\alpha_s}(x)\). در این صورت
\[
\begin{aligned}
s_1\cdot\bigl(s_2\cdot(a\otimes b)\bigr)
  &=s_1\cdot\bigl((s_2a)\otimes b\bigr)\\
  &=(s_1(s_2a))\otimes b\\
  &=((s_1s_2)a)\otimes b=(s_1s_2)\cdot(a\otimes b).
\end{aligned}
\]
سایر اصول مدولی نیز به همین صورت نتیجه می‌شوند.
\end{frame}

% Source: sources/lectures/20 tensor product 2.pdf | PDF page: 3 | item: S-linearity of f tensor g
\begin{frame}[t]{همریختی القایی و عمل چپ}
\label{lec20:left-tensor-morphism}
تابع \(f\otimes g\) قبلاً به عنوان همریختی گروهی تعریف شده است. برای نشان دادن \(S\)-خطی بودن آن داریم
\[
\begin{aligned}
(f\otimes g)\bigl(s\cdot(a\otimes b)\bigr)
  &=(f\otimes g)((sa)\otimes b)\\
  &=f(sa)\otimes g(b)\\
  &=(sf(a))\otimes g(b)\\
  &=s\cdot\bigl(f(a)\otimes g(b)\bigr)\\
  &=s\cdot(f\otimes g)(a\otimes b).
\end{aligned}
\]
\hfill\(\qed\)
\end{frame}

% Source: sources/lectures/20 tensor product 2.pdf | PDF page: 4 | item: dual right S action and tensor morphism theorem; proof referred to preceding theorem
\begin{frame}[t]{ساختار مدول راست روی ضرب تانسوری}
\label{lec20:right-tensor-action}
\begin{theorem}
فرض کنید \(C_R\)، \({}_R D_S\)، \(C'_R\) و \({}_R D'_S\) مدول باشند.
\begin{enumerate}
\item \(C\otimes_R D\) با عمل
\[
    (c\otimes d)\cdot s=c\otimes(ds)
\]
یک \(S\)-مدول راست است.
\item اگر \(h:C\to C'\) یک همریختی \(R\)-مدولی و
\(k:D\to D'\) یک همریختی دومدول‌های \(R\)-\(S\) باشد، آنگاه
\[
    h\otimes k:C\otimes_R D\longrightarrow C'\otimes_R D'
\]
یک همریختی \(S\)-مدولی است.
\end{enumerate}
\end{theorem}
\textbf{اثبات.} کاملاً مشابه قضیهٔ قبل است.
\end{frame}

% Source: sources/lectures/20 tensor product 2.pdf | PDF page: 5 | item: bilinear map over a commutative ring
\begin{frame}[t]{تابع دوخطی روی حلقهٔ جابجایی}
\label{lec20:bilinear-map}
\begin{definition}
فرض کنید \(R\) حلقه‌ای جابجایی و \(A,B,C\) سه \(R\)-مدول باشند.
تابع \(f:A\times B\to C\) را \textbf{دوخطی} می‌گوییم هرگاه
\[
    f(a_1+a_2,b)=f(a_1,b)+f(a_2,b),
\]
\[
    f(a,b_1+b_2)=f(a,b_1)+f(a,b_2),
\]
و برای هر \(r\in R\)،
\[
    f(ar,b)=r f(a,b)=f(a,rb).
\]
\end{definition}
\end{frame}

% Source: sources/lectures/20 tensor product 2.pdf | PDF page: 5 | item: R-linear universal property for a bilinear map
\begin{frame}[t]{خاصیت جهانی برای تابع دوخطی}
\label{lec20:linear-universal-property}
\begin{theorem}
اگر \(R\) جابجایی و \(f:A\times B\to C\) دوخطی باشد، یک همریختی یکتای \(R\)-مدولی
\[
    \bar f:A\otimes_R B\longrightarrow C
\]
وجود دارد که \(\bar f(a\otimes b)=f(a,b)\).
\end{theorem}
\textbf{اثبات.} خاصیت جهانیِ گروهی، \(\bar f\) را می‌دهد. چون \(R\) جابجایی است، ساختارهای مدولیِ لازم را داریم و
\[
\begin{aligned}
\bar f\bigl(r\cdot(a\otimes b)\bigr)
 &=\bar f(ra\otimes b)=f(ra,b)\\
 &=f(ar,b)=r f(a,b)=r\bar f(a\otimes b).
\end{aligned}
\]
\hfill\(\qed\)
\end{frame}

% Source: sources/lectures/20 tensor product 2.pdf | PDF page: 6 | item: unit isomorphisms and balanced multiplication map
\begin{frame}[t]{ضرب تانسوری با حلقه}
\label{lec20:unit-isomorphisms}
\begin{theorem}
برای مدول‌های \(A_R\) و \({}_R B\)، یکریختی‌های مدولی زیر برقرارند:
\[
    R\otimes_R B\cong B,\qquad A\otimes_R R\cong A.
\]
\end{theorem}
\textbf{اثباتِ یکریختی اول.} بگذارید
\[
    f:R\times B\longrightarrow B,\qquad f(r,b)=rb.
\]
این تابع در هر متغیر جمع‌پذیر است و
\[
    f(rr',b)=(rr')b=r(r'b)=f(r,r'b).
\]
پس \(f\) خطی میانی است و همریختی گروهیِ
\(\bar f:R\otimes_R B\to B\)، با \(\bar f(r\otimes b)=rb\)، وجود دارد.
\end{frame}

% Source: sources/lectures/20 tensor product 2.pdf | PDF page: 6 | item: linearity of multiplication-induced map
\begin{frame}[t]{خطی بودن نگاشتِ ضرب}
\label{lec20:unit-map-linear}
برای \(r'\in R\) داریم
\[
\begin{aligned}
\bar f\bigl(r'\cdot(r\otimes b)\bigr)
  &=\bar f(r'r\otimes b)\\
  &=f(r'r,b)=(r'r)b\\
  &=r'(rb)=r'f(r,b)\\
  &=r'\bar f(r\otimes b).
\end{aligned}
\]
بنابراین \(\bar f\) یک همریختی \(R\)-مدولی است. برای اثبات یکریختی بودن، وارون آن را تعریف می‌کنیم.
\end{frame}

% Source: sources/lectures/20 tensor product 2.pdf | PDF page: 7 | item: inverse g and compositions; second unit isomorphism similar
\begin{frame}[t]{وارون نگاشتِ ضرب}
\label{lec20:unit-inverse}
بگذارید \(g:B\to R\otimes_R B\)، \(g(b)=1\otimes b\). آنگاه
\[
    g(b_1+b_2)=1\otimes(b_1+b_2)=g(b_1)+g(b_2),
\]
\[
    g(rb)=1\otimes rb=r\otimes b=r(1\otimes b)=rg(b).
\]
همچنین
\[
    \bar f g(b)=\bar f(1\otimes b)=b,
\]
\[
    g\bar f(r\otimes b)=g(rb)=1\otimes rb=r\otimes b.
\]
پس \(\bar f\) و \(g\) وارون یکدیگرند. اثبات یکریختی دوم دقیقاً مشابه است. \hfill\(\qed\)
\end{frame}

% Source: sources/lectures/20 tensor product 2.pdf | PDF page: 7 | item: comparison of Q tensor Q over Z and over Q
\begin{frame}[t]{دو حلقهٔ پایه برای یک مثال}
\label{lec20:rational-examples}
\begin{example}
قبلاً دیدیم که به عنوان \(\mathbb Z\)-مدول،
\[
    \mathbb Q\otimes_{\mathbb Z}\mathbb Q\cong\mathbb Q.
\]
از قضیهٔ ضرب تانسوری با حلقه نیز، به عنوان \(\mathbb Q\)-مدول، داریم
\[
    \mathbb Q\otimes_{\mathbb Q}\mathbb Q\cong\mathbb Q.
\]
\end{example}
\end{frame}

% Source: sources/lectures/20 tensor product 2.pdf | PDF page: 8 | item: tensoring an exact sequence need not preserve initial zero
\begin{frame}[t]{ضرب تانسوری و صفرِ ابتدای رشته}
\label{lec20:exactness-example}
رشتهٔ دقیق زیر را در نظر بگیرید:
\[
    0\longrightarrow\mathbb Z\xrightarrow{\ f\ }\mathbb Z
    \xrightarrow{\ g\ }\mathbb Z_2\longrightarrow0,
\qquad f(a)=2a,\quad g(n)=\bar n.
\]
آیا پس از ضرب تانسوری با \(\mathbb Z_2\)، رشتهٔ زیر دقیق است؟
\[
\begin{gathered}
0\longrightarrow\mathbb Z\otimes_{\mathbb Z}\mathbb Z_2
\xrightarrow{\ f\otimes1\ }\mathbb Z\otimes_{\mathbb Z}\mathbb Z_2\\
\xrightarrow{\ g\otimes1\ }\mathbb Z_2\otimes_{\mathbb Z}\mathbb Z_2
\longrightarrow0.
\end{gathered}
\]
پاسخ منفی است؛ صفرِ سمت چپ لزوماً حفظ نمی‌شود.
\end{frame}

% Source: sources/lectures/20 tensor product 2.pdf | PDF page: 8 | item: vanishing of f tensor identity and failure of injectivity
\begin{frame}[t]{از بین رفتن یک‌به‌یکی بودن}
\label{lec20:noninjective-tensor-map}
برای \(a\in\mathbb Z\) و \(b\in\mathbb Z_2\)،
\[
\begin{aligned}
(f\otimes1)(a\otimes b)
 &=f(a)\otimes b=(2a)\otimes b\\
 &=a\otimes(2b)=a\otimes0=0.
\end{aligned}
\]
پس \(\operatorname{Im}(f\otimes1)=\{0\}\). از سوی دیگر
\[
    \mathbb Z\otimes_{\mathbb Z}\mathbb Z_2\cong\mathbb Z_2.
\]
بنابراین \(f\otimes1\) یک‌به‌یک نیست و صفرِ ابتدای رشته حفظ نشده است.
\end{frame}

% Source: sources/lectures/20 tensor product 2.pdf | PDF page: 8 | item: quotient tensor theorem for a right ideal and a left module
\begin{frame}[t]{ضرب تانسوریِ خارج‌قسمتِ حلقه}
\label{lec20:quotient-tensor}
\begin{theorem}
فرض کنید \(R\) یک حلقه، \(I\) یک ایده‌آل راستِ \(R\) و \(M\) یک \(R\)-مدول چپ باشد. آنگاه به عنوان \(\mathbb Z\)-مدول،
\[
    (R/I)\otimes_R M\cong M/IM.
\]
اگر \(R\) جابجایی باشد، این یکریختی، \(R\)-مدولی است.
\end{theorem}
\textbf{اثبات.} تابع زیر را تعریف می‌کنیم:
\[
    f:(R/I)\times M\longrightarrow M/IM,
    \qquad f(r+I,m)=rm+IM.
\]
\end{frame}

% Source: sources/lectures/20 tensor product 2.pdf | PDF page: 9 | item: well-definedness and source's balancedness verification prompt
\begin{frame}[t]{خوش‌تعریف بودن نگاشتِ خارج‌قسمت}
\label{lec20:quotient-well-defined}
اگر \(r+I=r'+I\)، آنگاه \(r=r'+a\) برای یک \(a\in I\). بنابراین
\[
\begin{aligned}
    rm+IM&=(r'+a)m+IM\\
          &=r'm+am+IM=r'm+IM.
\end{aligned}
\]
پس \(f\) خوش‌تعریف است. افزون بر این، \(f\) تابعی خطی میانی است. \alert{چرا؟}
\medskip

بنا بر خاصیت جهانی، همریختی \(\mathbb Z\)-مدولیِ زیر وجود دارد:
\[
    \bar f:(R/I)\otimes_R M\longrightarrow M/IM,
    \quad\bar f((r+I)\otimes m)=rm+IM.
\]
\end{frame}

% Source: sources/lectures/20 tensor product 2.pdf | PDF page: 9 | item: surjectivity and kernel calculation setup
\begin{frame}[t]{پوشایی و آغاز بررسی هسته}
\label{lec20:quotient-kernel-setup}
نگاشت \(\bar f\) پوشاست؛ زیرا برای هر \(m\in M\)،
\[
    \bar f((1+I)\otimes m)=m+IM.
\]
اکنون اگر
\[
    \bar f\left(\sum_{i=1}^n(r_i+I)\otimes m_i\right)=0,
\]
آنگاه
\[
    \sum_{i=1}^n r_i m_i\in IM.
\]
پس می‌توان نوشت
\[
    \sum_{i=1}^n r_i m_i=\sum_{t=1}^k a_t m_t,\qquad a_t\in I.
\]
\end{frame}

% Source: sources/lectures/20 tensor product 2.pdf | PDF pages: 9-10 | item: injectivity of quotient tensor map
\begin{frame}[t]{صفر بودن هسته}
\label{lec20:quotient-injective}
با استفاده از رابطهٔ خطی میانی و جمع‌پذیری داریم
\[
\begin{aligned}
\sum_{i=1}^n(r_i+I)\otimes m_i
  &=\sum_{i=1}^n(1+I)\otimes r_i m_i\\
  &=(1+I)\otimes\left(\sum_{i=1}^n r_i m_i\right)\\
  &=(1+I)\otimes\left(\sum_{t=1}^k a_t m_t\right)\\
  &=\sum_{t=1}^k(a_t+I)\otimes m_t=0.
\end{aligned}
\]
پس \(\bar f\) یک‌به‌یک و در نتیجه یک یکریختی گروهی است.
\end{frame}

% Source: sources/lectures/20 tensor product 2.pdf | PDF page: 10 | item: R-linearity when R is commutative
\begin{frame}[t]{حالت حلقهٔ جابجایی}
\label{lec20:quotient-linear}
اگر \(R\) جابجایی باشد، \(R/I\) یک \(R\)-\(R\) دومدول است و عمل اسکالر روی ضرب تانسوری به صورت زیر است:
\[
    r'\cdot((r+I)\otimes m)=(r'r+I)\otimes m.
\]
در این صورت
\[
\begin{aligned}
\bar f\bigl(r'\cdot((r+I)\otimes m)\bigr)
  &=\bar f((r'r+I)\otimes m)\\
  &=r'rm+IM=r'(rm+IM)\\
  &=r'\bar f((r+I)\otimes m).
\end{aligned}
\]
پس یکریختیِ به‌دست‌آمده، \(R\)-مدولی است. \hfill\(\qed\)
\end{frame}

% Source: sources/lectures/20 tensor product 2.pdf | PDF page: 10 | item: tensor product of finite cyclic groups and quotient-chain proof
\begin{frame}[t]{ضرب تانسوریِ دو گروه دوری}
\label{lec20:cyclic-tensor}
\begin{corollary}
اگر \(d=(m,n)\)، آنگاه
\[
    \mathbb Z_m\otimes_{\mathbb Z}\mathbb Z_n\cong\mathbb Z_d.
\]
\end{corollary}
\textbf{اثبات.}
\[
\begin{aligned}
\mathbb Z_m\otimes_{\mathbb Z}\mathbb Z_n
 &\cong(\mathbb Z/m\mathbb Z)\otimes_{\mathbb Z}\mathbb Z_n
 \cong\mathbb Z_n/m\mathbb Z_n\\
 &\cong\frac{\mathbb Z/n\mathbb Z}{(m\mathbb Z+n\mathbb Z)/n\mathbb Z}\\
 &\cong\mathbb Z/d\mathbb Z\cong\mathbb Z_d.
\end{aligned}
\]
\hfill\(\qed\)
\end{frame}

% Source: sources/lectures/20 tensor product 2.pdf | PDF page: 11 | item: hypotheses and statement of associativity
\begin{frame}[t]{شرکت‌پذیریِ ضرب تانسوری}
\label{lec20:associativity}
فرض کنید \(R,S\) دو حلقه و
\[
    A_R,\qquad {}_R B_S,\qquad {}_S C
\]
مدول باشند. ساختارهای القاییِ
\[
    {}_R(B\otimes_S C),\qquad (A\otimes_R B)_S
\]
دو ضرب تانسوریِ زیر را ممکن می‌کنند.
\begin{theorem}
به عنوان \(\mathbb Z\)-مدول،
\[
    (A\otimes_R B)\otimes_S C\cong A\otimes_R(B\otimes_S C).
\]
\end{theorem}
\end{frame}

% Source: sources/lectures/20 tensor product 2.pdf | PDF page: 11 | item: fixed-c map T_c and both additivity checks
\begin{frame}[t]{نگاشتِ وابسته به یک عنصر ثابت}
\label{lec20:fixed-c-map}
\textbf{اثبات.} برای \(c\in C\)، تعریف می‌کنیم
\[
    T_c:A\times B\longrightarrow A\otimes_R(B\otimes_S C),
    \qquad T_c(a,b)=a\otimes(b\otimes c).
\]
این تابع در هر متغیر جمع‌پذیر است:
\[
\begin{aligned}
T_c(a_1+a_2,b)
 &=(a_1+a_2)\otimes(b\otimes c)\\
 &=T_c(a_1,b)+T_c(a_2,b),\\[2pt]
T_c(a,b_1+b_2)
 &=a\otimes(b_1\otimes c+b_2\otimes c)\\
 &=T_c(a,b_1)+T_c(a,b_2).
\end{aligned}
\]
\end{frame}

% Source: sources/lectures/20 tensor product 2.pdf | PDF page: 12 | item: R-balancing of T_c and induced group homomorphism
\begin{frame}[t]{خطی میانی بودن نگاشتِ نخست}
\label{lec20:fixed-c-balanced}
برای \(r\in R\)،
\[
\begin{aligned}
T_c(ar,b)
  &=ar\otimes(b\otimes c)\\
  &=a\otimes\bigl(r(b\otimes c)\bigr)\\
  &=a\otimes((rb)\otimes c)=T_c(a,rb).
\end{aligned}
\]
بنابراین \(T_c\) خطی میانی است و همریختی گروهیِ زیر وجود دارد:
\[
    \overline{T_c}:A\otimes_R B\longrightarrow A\otimes_R(B\otimes_S C),
\]
\[
    \overline{T_c}(a\otimes b)=a\otimes(b\otimes c).
\]
\end{frame}

% Source: sources/lectures/20 tensor product 2.pdf | PDF page: 12 | item: definition of f through T_c-bar
\begin{frame}[t]{نگاشتِ مرحلهٔ دوم}
\label{lec20:associativity-second-map}
حال تابع زیر را تعریف می‌کنیم:
\[
    f:(A\otimes_R B)\times C\longrightarrow A\otimes_R(B\otimes_S C),
\]
\[
\begin{aligned}
f\left(\sum_{i=1}^n a_i\otimes b_i,c\right)
  &:=\overline{T_c}\left(\sum_{i=1}^n a_i\otimes b_i\right)\\
  &=\sum_{i=1}^n a_i\otimes(b_i\otimes c).
\end{aligned}
\]
این تابع نیز خطی میانی است. \alert{چرا؟}
در ادامه سه شرط آن را بررسی می‌کنیم.
\end{frame}

% Source: sources/lectures/20 tensor product 2.pdf | PDF page: 12 | item: additivity of f in C and in A tensor B
\begin{frame}[t]{دو شرط جمع‌پذیری}
\label{lec20:associativity-additivity}
اگر \(x=\sum_{i=1}^n a_i\otimes b_i\)، آنگاه
\[
\begin{aligned}
f(x,c_1+c_2)
 &=\sum_{i=1}^n a_i\otimes\bigl(b_i\otimes(c_1+c_2)\bigr)\\
 &=\sum_{i=1}^n a_i\otimes(b_i\otimes c_1)
   +\sum_{i=1}^n a_i\otimes(b_i\otimes c_2)\\
 &=f(x,c_1)+f(x,c_2).
\end{aligned}
\]
همچنین، چون \(\overline{T_c}\) همریختی گروهی است،
\[
\begin{aligned}
f(x+x',c)&=\overline{T_c}(x+x')\\
 &=\overline{T_c}(x)+\overline{T_c}(x')=f(x,c)+f(x',c).
\end{aligned}
\]
\end{frame}

% Source: sources/lectures/20 tensor product 2.pdf | PDF page: 13 | item: S-balancing of the second map
\begin{frame}[t]{شرط خطی میانی نسبت به حلقهٔ دوم}
\label{lec20:associativity-s-balance}
برای \(s\in S\) و \(x=\sum_{i=1}^n a_i\otimes b_i\)،
\[
\begin{aligned}
f(xs,c)
 &=f\left(\sum_{i=1}^n a_i\otimes(b_is),c\right)\\
 &=\sum_{i=1}^n a_i\otimes((b_is)\otimes c)\\
 &=\sum_{i=1}^n a_i\otimes(b_i\otimes(sc))\\
 &=f(x,sc).
\end{aligned}
\]
پس \(f\) تابعی خطی میانی است و می‌توان بار دیگر خاصیت جهانی را به کار برد.
\end{frame}

% Source: sources/lectures/20 tensor product 2.pdf | PDF page: 13 | item: induced associator, similarly constructed inverse, and identity compositions
\begin{frame}[t]{یکریختیِ شرکت‌پذیری و وارون آن}
\label{lec20:associator-inverse}
همریختی گروهیِ حاصل چنین است:
\[
    \bar f:(A\otimes_R B)\otimes_S C\longrightarrow A\otimes_R(B\otimes_S C),
\]
\[
    \bar f((a\otimes b)\otimes c)=a\otimes(b\otimes c).
\]
با استدلالی کاملاً مشابه، همریختی \(\bar g\) در جهت عکس وجود دارد که
\[
    \bar g(a\otimes(b\otimes c))=(a\otimes b)\otimes c.
\]
در نتیجه
\[
    \bar g\bar f=1,\qquad\bar f\bar g=1.
\]
پس \(\bar f\) یک یکریختی \(\mathbb Z\)-مدولی است. \hfill\(\qed\)
\end{frame}
